\section{Mellin positivity and one-sided pole detection}
\label{app:mellin-positivity}

The real pole at \(s=1/2\) gives a negative bias in the smoothed
Liouville sum. It is natural to test whether one can prove an eventual
sign and thereby exclude right-hand zeros. This appendix proves
the relevant implications, including the positivity principle,
without assuming RH. An eventual sign is not proved here.

\begin{lemma}[The positivity principle]\label{lem:landau-positive}
Let \(f:[0,\infty)\to[0,\infty)\) be locally integrable, and suppose
its Laplace integral
\[
 \mathcal L(s)=\int_0^\infty f(t)e^{-st}\dd t
\]
converges for all sufficiently large real \(s\). If its abscissa
of convergence \(a\) is finite, then \(\mathcal L\) cannot extend
holomorphically to a neighborhood of the real point \(a\).
\end{lemma}
\begin{proof}
Suppose there is such a continuation, holomorphic in
\(\operatorname{Re}s>a\) and in a disc of radius \(r\) about \(a\).
Choose \(b=a+r/4\). The disc about \(b\) of radius \(r/2\)
lies in the union of those two domains. Differentiation in the
convergent integral gives the Taylor coefficients of
\(\mathcal L(b-z)\):
\[
 \frac1{j!}\int_0^\infty t^j f(t)e^{-bt}\dd t\ge0.
\]
The differentiations are justified by convergence at a real
point strictly between \(a\) and \(b\). For a real
\(z\in(r/4,r/2)\), the Taylor series converges. By Tonelli its
sum is
\[
 \int_0^\infty f(t)e^{-bt}
                 \sum_{j\ge0}\frac{(zt)^j}{j!}\dd t
 =\int_0^\infty f(t)e^{-(b-z)t}\dd t.
\]
But \(b-z<a\), contradicting the definition of the abscissa.
\end{proof}

The same principle applies to a nonnegative function \(A(X)\) on
\([X_0,\infty)\) with Mellin integral
\(\int_{X_0}^\infty A(X)X^{-s}\dd X/X\), by putting
\(X=X_0e^t\). Changing \(X_0\) only changes the meromorphic
continuation by an entire function.

\begin{theorem}[Both signs forced by a right-hand zero]
\label{thm:two-sided-smoothed}
Let \(\rho=\beta+i\gamma\) be a nontrivial zeta zero of order \(r\)
with \(\beta>1/2\). Define
\[
 C_\rho=\frac{r|\Gamma(\rho)\zeta(2\rho)|}
                     {|\zeta^{(r)}(\rho)|}>0.
\]
Then
\begin{align*}
 \limsup_{X\to\infty}
  \frac{B(X)}{X^\beta(\log X)^{r-1}}&\ge C_\rho,\\
 \liminf_{X\to\infty}
  \frac{B(X)}{X^\beta(\log X)^{r-1}}&\le-C_\rho.
\end{align*}
\end{theorem}
\begin{proof}
Put \(\Phi(s)=\Gamma(s)\zeta(2s)/\zeta(s)\).
It is the Mellin transform in \eqref{eq:gamma} for
\(\operatorname{Re}s>1\). It has no singularity at any real
point greater than \(1/2\). For \(1/2<s<1\), this follows from
\(\zeta(s)=\eta(s)/(1-2^{1-s})<0\), where the alternating
series \(\eta(s)>0\); at one the pole of \(\zeta\) gives a
removable zero in the quotient. At \(\rho\) the leading
Laurent coefficient is
\[
 A_\rho=\frac{r!\Gamma(\rho)\zeta(2\rho)}
                     {\zeta^{(r)}(\rho)}.
\]
The numerator is nonzero since \(\operatorname{Re}(2\rho)>1\).

Suppose the first inequality fails. Choose
\(0<c<C_\rho\) such that
\[
 A(X)=cX^\beta(\log X)^{r-1}-B(X)\ge0
\]
for every \(X\ge X_0\), with \(X_0>1\). Its Mellin transform
converges for real \(s>1\) and continues meromorphically as
\[
 \frac{c(r-1)!}{(s-\beta)^r}-\Phi(s)+E(s),
\]
where \(E\) is entire, accounting for the finite initial interval.
The real pole at \(\beta\) forces the abscissa of the positive
integral to be at least \(\beta\). If that abscissa were greater
than \(\beta\), the displayed continuation would be holomorphic
there, contradicting Lemma~\ref{lem:landau-positive}. Hence it
is exactly \(\beta\).

Positivity now gives, for \(\delta>0\),
\[
 \left|\int_{X_0}^\infty A(X)X^{-\beta-\delta-i\gamma}
                       \frac{\dd X}{X}\right|
 \le\int_{X_0}^\infty A(X)X^{-\beta-\delta}
                       \frac{\dd X}{X}.
\]
Multiply by \(\delta^r\) and let \(\delta\downarrow0\).
The complex pole on the left gives \(|A_\rho|\); the real
pole on the right gives \(c(r-1)!\). Therefore
\(c\ge |A_\rho|/(r-1)!=C_\rho\), a contradiction.
For the second inequality apply the identical argument to
\(A(X)=cX^\beta(\log X)^{r-1}+B(X)\).
\end{proof}

This strengthens the absolute-value pole test of
Proposition~\ref{prop:two-pole-bounds} to both one-sided
oscillations. It still gives lower bounds, not the square-root
upper bound needed to exclude \(\rho\).

\begin{proposition}[The real pole and an eventual sign]
\label{prop:eventual-smoothed-sign}
Put
\[
 c_0=-\frac{\sqrt\pi}{2\zeta(1/2)}>0.
\]
Unconditionally,
\[
 \liminf_{X\to\infty}\frac{B(X)}{\sqrt X}\le-c_0.
\]
If in addition \(B(X)\le0\) for all sufficiently large \(X\),
then RH holds, every nontrivial zero is simple, and each such
zero satisfies
\[
 \left|\frac{\Gamma(\rho)\zeta(2\rho)}{\zeta'(\rho)}\right|
 \le c_0.
\]
\end{proposition}
\begin{proof}
The function \(\Phi\) has a simple real pole at \(1/2\)
with residue \(-c_0\). Suppose the asserted limit inferior
were greater than \(-c_0\). Choose \(0<c<c_0\) such that
\(B(X)+c\sqrt X\ge0\) eventually. The Mellin transform of
this nonnegative function has continuation
\[
 \Phi(s)+\frac c{s-1/2}+E(s).
\]
Its pole at \(1/2\) has negative residue \(c-c_0\).
The positivity principle and absence of real singularities to
the right force its abscissa to be \(1/2\). The integral is
nonnegative on the real axis there, but the continuation
tends to \(-\infty\) as \(s\downarrow1/2\), a contradiction.

Now assume \(B\le0\) eventually, and apply the positivity
principle to \(-B\). Its Mellin continuation is
\(-\Phi+E\), with positive residue \(c_0\) at \(1/2\).
Its abscissa is exactly \(1/2\), by the same real-axis
argument. Thus \(\Phi\) is holomorphic in
\(\operatorname{Re}s>1/2\), which excludes all right-hand
zeros. The functional equation gives RH.

At a critical zero of order \(r\), the numerator of \(\Phi\)
is nonzero by nonvanishing of \(\zeta\) on the line
\(\operatorname{Re}s=1\). Positivity bounds the complex
Mellin integral at \(s=\rho+\delta\) by its real value at
\(s=1/2+\delta\), the latter being \(c_0/\delta+O(1)\).
The complex pole has order \(r\), so \(r\le1\), and comparison
of residues gives the displayed inequality.
\end{proof}

An eventual-sign argument would therefore need more than the
square-root estimate with an arbitrary \(X^\varepsilon\) loss:
it would also establish simplicity. Neither the sign assertion
nor the required upper bound for \(B\) has been proved.
The negative-bias lower bound alone does not control positive
excursions or rule out right-hand zeros.

\begin{proposition}[The real-pole normalization of local energy]
\label{prop:local-energy-lower}
With \(c_0\) as above, unconditionally
\[
 \limsup_{X\to\infty}\frac{\mathcal H_B(X)}X
 \ge\frac32c_0^2=\frac{3\pi}{8\zeta(1/2)^2}.
\]
\end{proposition}
\begin{proof}
There is nothing to prove if the limit superior is infinite.
Otherwise choose \(C\) greater than it, so that
\(\mathcal H_B(X)\le CX\) for all sufficiently large \(X\).
For \(\eta>0\), Tonelli and the substitution \(Y=tX\) give
\[
 \int_0^\infty\mathcal H_B(X)X^{-1-\eta}\frac{\dd X}{X}
 =r(1+\eta)P(\eta),
\]
where
\[
 r(s)=\int_{1/2}^2t^s\frac{\dd t}{t}
     =\frac{2^s-2^{-s}}s,\qquad
 P(\eta)=\int_0^\infty|B(Y)|^2Y^{-1-\eta}\frac{\dd Y}{Y}.
\]
The small-\(X\) part is uniformly finite near \(\eta=0\),
by exponential decay. The assumed large-\(X\) upper bound gives
\[
 P(\eta)\le\frac{C/\eta+O(1)}{r(1+\eta)}.
\]
For \(\delta>0\), Cauchy--Schwarz on \([1,\infty)\) now gives
\[
 \left|\int_1^\infty B(Y)Y^{-1/2-\delta}
                       \frac{\dd Y}{Y}\right|^2
 \le\frac{P(\delta)}{\delta}.
\]
The absolute integral is convergent by the same estimate.
It agrees with \(\Phi(1/2+\delta)\) up to the entire
small-\(Y\) term, by continuation from
\(\operatorname{Re}s>1\). Multiplying by \(\delta^2\) and
letting \(\delta\downarrow0\) gives
\(c_0^2\le C/r(1)=2C/3\).
Let \(C\) decrease to the limit superior.
\end{proof}

For the comparison function \(B_0(X)=-c_0\sqrt X\), the local energy
is exactly \((3/2)c_0^2X\). Thus the constant in this pole comparison
is the one dictated by the real residue. This proposition supplies
a lower constraint on any proposed upper estimate; it does not
establish such an estimate for the actual sum.

\section{Elementary prime bounds and exact checks}
\label{app:attack-checks}

Here are details of the prime estimate used in
Theorem~\ref{thm:finite-prime}. If
\(\vartheta(t)=\sum_{p\le t}\log p\), then
\(\prod_{n<p\le2n}p\) divides \(\binom{2n}{n}\le4^n\).
Summing over dyadic intervals gives
\(\vartheta(t)\le4(\log2)t\) for \(t\ge2\). Separating
primes at \(\sqrt t\) gives
\[
 \pi(t)\le\sqrt t+\frac{2\vartheta(t)}{\log t}
        \le7\,\frac t{\log t}.
\]
The last inequality uses \(\log t\le\sqrt t\).
Partial summation gives
\[
 \sum_{p\le y}p^{-1/2}
 =\frac{\pi(y)}{\sqrt y}
  +\frac12\int_2^y\frac{\pi(t)}{t^{3/2}}\dd t
 \ll\frac{\sqrt y}{\log y}.
\]
For instance split the integral at \(\sqrt y\); the first
part is \(O(y^{1/4})\) and on the second part
\(\log t\ge\frac12\log y\). This proves the asserted order
without the prime number theorem or any zero-free input.
Similarly
\(\sum_{p\le y}(\log p)/p\ll\log y\), by partial summation
of \(\vartheta\), proving the derivative bound used in
\eqref{eq:prime-tail-double-zero}.

The new checker \path{checks/attack_exact.py} uses rational
arithmetic for the following independent finite tests:
the Legendre reproducing kernel on polynomial monomials,
its evaluation norm bound, the positive coefficients of the
diagonal error, both finite-prime convolution identities,
the coprime-pair decomposition, and the prime-dilation Gram
identity. It also checks an exponential interval certificate
for \(B(1)>0\) and \(B(3)<-1/20\).
The latter rules out assigning one sign to every scale
product in \eqref{eq:lambda-log-renewal}; it is not a test
of an eventual sign.

The checker proves none of the infinite arithmetic estimates.
The analytic arguments are supplied in the text and have
not been formalized in Lean. The original finite Gram
formalization retains exactly the scope stated in the
reproduction appendix. No new RH proof, zero-free
half-plane or novelty priority is asserted in this revision.
